\subsection{Definitions and statement}

Work over \(\C\), and let
\[
  V=\Span\{x_{ij}:0\leq i,j\leq3\},\qquad \dim V=16.
\]

\begin{definition}
The Chow rank of \(f\in\Sym^4V\) is the least integer \(r\) for which
\[
  f=\sum_{s=1}^{r}\ell_{s1}\ell_{s2}\ell_{s3}\ell_{s4},
  \qquad \ell_{st}\in V.
\]
Each product \(\ell_1\ell_2\ell_3\ell_4\) is a Chow term.
\end{definition}

\begin{theorem}\label{n4:thm:main}
\[
  \boxed{\ChowRank(\perm_4)=8}.
\]
\end{theorem}

Section~\ref{n4:sec:upper} gives an eight-term decomposition. The preceding
sections prove that no decomposition with seven terms exists.

\subsection{Second derivatives and the Koszul flattening}
\label{n4:sec:koszul}

For \(f\in\Sym^4V\), let
\[
  C_{2,2}(f):\Sym^2V^*\longrightarrow\Sym^2V,
  \qquad D\longmapsto D(f),
\]
and set \(E_f:=\im C_{2,2}(f)\). Define
\[
  \delta:\Sym^2V\otimes V\longrightarrow V\otimes\bigwedge\nolimits^2V
\]
by
\begin{equation}\label{n4:eq:delta}
  \delta\bigl((uv)\otimes w\bigr)
  =v\otimes(u\wedge w)+u\otimes(v\wedge w),
\end{equation}
extended linearly. We use the flattening
\[
  K(f):=\delta\circ\bigl(C_{2,2}(f)\otimes\id_V\bigr),
\]
for which
\begin{equation}\label{n4:eq:image-k}
  \im K(f)=\delta(E_f\otimes V).
\end{equation}
In monomial bases the middle catalecticant is symmetric. Transposing gives
\begin{equation}\label{n4:eq:row-containment}
  \row K(f)\subseteq E_f\otimes V^*.
\end{equation}

Write \(E:=E_{\perm_4}\). Direct differentiation gives
\begin{equation}\label{n4:eq:E-basis}
  E=\Span\left\{
    x_{ia}x_{jb}+x_{ib}x_{ja}:
    0\leq i<j\leq3,\ 0\leq a<b\leq3
  \right\}.
\end{equation}
These quadrics have distinct row-pair and column-pair supports, so
\(\dim E=\binom42^2=36\).

\begin{proposition}\label{n4:prop:base-ranks}
\[
  \rank K(\perm_4)=560.
\]
For every Chow term \(T=\ell_1\ell_2\ell_3\ell_4\),
\[
  \rank K(T)\leq92,
\]
with equality when \(\ell_1,\ldots,\ell_4\) are linearly independent.
\end{proposition}

\begin{proof}
The domain of
\(\delta|_{E\otimes V}\) has dimension \(36\cdot16=576\). If all first
derivatives of \(g\in\Sym^3V\) lie in \(E\), define
\[
  \iota(g):=\sum_{a=0}^{15}
  \frac{\partial g}{\partial x_a}\otimes x_a.
\]
Equation~\eqref{n4:eq:delta} gives \(\delta(\iota(g))=0\), while tensor
multiplication sends \(\iota(g)\) to \(3g\) by Euler's identity. Thus
\(\iota\) is injective.

For every three-row set \(I\subset\{0,1,2,3\}\) and three-column set
\(J\subset\{0,1,2,3\}\), let
\[
  p_{I,J}:=\perm(x_{ij})_{i\in I,j\in J}.
\]
The sixteen such cubics have disjoint supports, and all their first
derivatives lie in \(E\). They therefore give sixteen independent kernel
vectors. Hence \(\rank K(\perm_4)\leq576-16=560\). In the bases from
\eqref{n4:eq:E-basis}, exact elimination finds a \(560\times560\) integer
minor that is nonzero modulo \(1000003\). The same determinant is a nonzero
integer, so the reverse inequality holds in characteristic zero.

For a term with independent factors, a \(GL(V)\)-change of coordinates gives
\(T=x_0x_1x_2x_3\). Then
\[
  E_T=\Span\{x_ix_j:0\leq i<j\leq3\},\qquad \dim E_T=6.
\]
The four squarefree cubics in \(x_0,x_1,x_2,x_3\) give four independent
kernel vectors, so \(\rank K(T)\leq96-4=92\). Exact elimination finds a
nonzero \(92\times92\) minor modulo \(1000003\), giving equality. Independent
factor quadruples form a dense open set. Since the locus where matrix rank
exceeds \(92\) is open, no degenerate quadruple can have larger rank: such an
open set would meet the dense independent locus. Thus every Chow term has
rank at most \(92\).
\end{proof}

\begin{remark}
The modular calculations certify nonzero integer minors; they are not
probabilistic rank estimates. Section~\ref{n4:sec:computation} gives the exact
reconstruction commands.
\end{remark}

\subsection{Low-rank quadrics in the derivative space}
\label{n4:sec:classification}

Index the basis in \eqref{n4:eq:E-basis} by the six row pairs and six column
pairs. A quadric \(q\in E\) corresponds to a \(6\times6\) coefficient matrix
\(C\). For a fixed row pair \(i<j\), write the corresponding row as a
symmetric zero-diagonal \(4\times4\) matrix
\[
  D_{ij}[a,b]=C[(i,j),(a,b)]\quad(a\neq b),\qquad D_{ij}[a,a]=0.
\]
Grouping the variables by matrix rows, the symmetric coefficient matrix is
\begin{equation}\label{n4:eq:block-matrix}
  M(q)=(D_{ij})_{0\leq i,j\leq3},\qquad
  D_{ji}=D_{ij},\quad D_{ii}=0.
\end{equation}
Using the Hessian instead changes this matrix only by a nonzero scalar.

\begin{lemma}\label{n4:lem:zero-diagonal}
Let \(D\neq0\) be a symmetric zero-diagonal matrix over \(\C\).
\begin{enumerate}[label=(\roman*)]
  \item \(\rank D\neq1\).
  \item If \(\rank D=2\), and \(H\) is symmetric and zero-diagonal with
  \(\im H\subseteq\im D\), then \(H\) is a scalar multiple of \(D\).
\end{enumerate}
\end{lemma}

\begin{proof}
A rank-one symmetric matrix has the form \(\lambda uu^{\mathsf T}\). A zero
diagonal gives \(\lambda u_i^2=0\) for every \(i\), contradicting
\(D\neq0\).

If \(\rank D=2\), choose independent \(a,b\) such that
\[
  D=ab^{\mathsf T}+ba^{\mathsf T}.
\]
The zero diagonal gives \(a_kb_k=0\) for every \(k\), so the supports of
\(a\) and \(b\) are disjoint. Symmetry and
\(\im H\subseteq\Span(a,b)\) give
\[
  H=[\,a\ b\,]R[\,a\ b\,]^{\mathsf T}
\]
for a symmetric \(2\times2\) matrix \(R\). The diagonal entries on the
nonempty supports of \(a\) and \(b\) force \(R_{11}=R_{22}=0\), leaving only
one off-diagonal parameter. Hence \(H\) is a scalar multiple of \(D\).
\end{proof}

\begin{theorem}\label{n4:thm:rank-four-classification}
If \(0\neq q\in E\) and \(\rank M(q)\leq4\), then
\[
  M(q)=U\otimes W,
\]
where \(U\) and \(W\) are symmetric zero-diagonal \(4\times4\) matrices of
rank \(2\). In particular, \(\rank M(q)=4\).
\end{theorem}

\begin{proof}
Choose a nonzero block \(D:=D_{ij}\). The associated principal submatrix is
\[
  N=\begin{bmatrix}0&D\\D&0\end{bmatrix},
  \qquad \rank N=2\rank D\leq\rank M(q)\leq4.
\]
Lemma~\ref{n4:lem:zero-diagonal} implies \(\rank D=2\), hence
\(\rank N=4=\rank M(q)\). Thus all remaining block columns lie in the column
space of \(N\). For each \(k\),
\[
  \im D_{ik},\ \im D_{jk}\subseteq\im D,
\]
so Lemma~\ref{n4:lem:zero-diagonal}(ii) gives scalars \(a_k,b_k\) with
\[
  D_{ik}=a_kD,\qquad D_{jk}=b_kD.
\]
Let \(G\) be a generalized inverse satisfying \(DGD=D\), and put
\[
  N^-=\begin{bmatrix}0&G\\G&0\end{bmatrix}.
\]
Because the full matrix and \(N\) have the same rank, the generalized Schur
complement vanishes. Therefore
\begin{align*}
  D_{kl}
  &=\begin{bmatrix}D_{ik}\\D_{jk}\end{bmatrix}^{\mathsf T}
    N^-
    \begin{bmatrix}D_{il}\\D_{jl}\end{bmatrix}\\
  &=(a_kb_l+b_ka_l)D.
\end{align*}
All blocks are scalar multiples of \(D\), so \(M(q)=U\otimes W\) with
\(W=D\) and \(U\) symmetric and zero-diagonal. Finally,
\[
  4=\rank M(q)=\rank U\cdot\rank W.
\]
Neither nonzero zero-diagonal factor can have rank one, so both ranks are two.
\end{proof}

\begin{corollary}\label{n4:cor:intersection}
For every subspace \(L\subseteq V\) with \(\dim L\leq4\),
\[
  \dim\bigl(E\cap\Sym^2L\bigr)\leq1.
\]
Consequently, if \(F=E_T\) for a Chow term \(T\), then
\[
  \dim(E\cap F)\leq1.
\]
\end{corollary}

\begin{proof}
If \(0\neq q\in E\cap\Sym^2L\), then \(\rank M(q)\leq\dim L\leq4\).
Theorem~\ref{n4:thm:rank-four-classification} shows that the rank is four and
that the essential variable space is \(\im U\otimes\im W\), also of dimension
four. Thus \(\dim L=4\) and \(L\) equals this essential space.

Tensor-product factors of nonzero two-dimensional subspaces are unique: if
\(S_1\otimes T_1=S_2\otimes T_2\), then \(S_1=S_2\) and \(T_1=T_2\).
Within a fixed two-dimensional image, Lemma~\ref{n4:lem:zero-diagonal}(ii)
makes the rank-two symmetric zero-diagonal matrix unique up to scale. Hence
any two nonzero elements of \(E\cap\Sym^2L\) are proportional.

For a Chow term, every second derivative uses only
\(L=\Span(\ell_1,\ldots,\ell_4)\), so \(F\subseteq\Sym^2L\). The second
claim follows, including repeated or dependent factors.
\end{proof}

\subsection{The chart calculation}\label{n4:sec:chart}

For \(0\neq v\in V\), define
\begin{equation}\label{n4:eq:psi}
  \psi_v:\Sym^2V/E\longrightarrow
  \coker\bigl(\delta(E\otimes V)\bigr),\qquad
  [q]\longmapsto[\delta(q\otimes v)].
\end{equation}
Its domain has dimension \(136-36=100\). Since
\(\delta(v^2\otimes v)=0\), and \(v^2\notin E\) by
Theorem~\ref{n4:thm:rank-four-classification},
\begin{equation}\label{n4:eq:obvious-kernel}
  0\neq[v^2]\in\ker\psi_v,qquad \rank\psi_v\leq99.
\end{equation}

\begin{proposition}\label{n4:prop:chart}
For every \(0\neq v\in V\),
\[
  \rank\psi_v=99,qquad \ker\psi_v=\Span([v^2]).
\]
\end{proposition}

\begin{proof}
First take the affine chart \(v_{00}\neq0\). Use the lexicographic bases
\[
  \Sym^2V:\ x_ax_b\ (0\leq a\leq b<16),\qquad
  V\otimes\bigwedge\nolimits^2V:\
  x_a\otimes(x_b\wedge x_c)\ (b<c).
\]
The program selects \(560\) independent columns of
\(\delta(E\otimes V)\), their \(560\) pivot rows, \(99\) classes in
\(\Sym^2V/E\), and \(99\) additional rows. For
\(v=\sum_{k=0}^{15}v_kx_k\), with \(v_0=v_{00}\), the resulting minor is
\begin{equation}\label{n4:eq:659}
  \mathcal M(v)=
  \begin{bmatrix}A_0&Q_0(v)\\A_1&Q_1(v)\end{bmatrix},
\end{equation}
where \(A_0\) is \(560\times560\) and the right side has \(99\) columns.
Exact rational elimination gives \(\det A_0=1\). Put
\[
  S(v):=Q_1(v)-A_1A_0^{-1}Q_0(v)=\sum_{k=0}^{15}v_kS_k.
\]
Direct calculation over \(\Q\) gives
\begin{equation}\label{n4:eq:S0}
  \det S_0=-32768.
\end{equation}
For \(k=1,\ldots,15\), set \(N_k=S_0^{-1}S_k\). The exact output verifies
that every \(N_k\) has zero diagonal, that their combined support has \(156\)
positions, and that the directed graph formed by these positions is acyclic.
A common topological ordering therefore makes all \(N_k\) strictly upper
triangular. Hence
\begin{align}
  \det\mathcal M(v)
  &=\det A_0\det S_0
    \det\left(v_0I_{99}+\sum_{k=1}^{15}v_kN_k\right)\notag\\
  &=-32768\,v_{00}^{99}.\label{n4:eq:chart-determinant}
\end{align}
This is nonzero when \(v_{00}\neq0\), so \(\rank\psi_v\geq99\).
Equation~\eqref{n4:eq:obvious-kernel} gives equality and identifies the kernel.

The group \(S_4\times S_4\) permutes the matrix rows and columns, preserves
\(\perm_4\), \(E\), and \(\delta\), and acts transitively on the sixteen
coordinates. Every nonzero \(v\) has a nonzero coordinate, which can be moved
to \((0,0)\). This reduces the general case to the chart already proved.
\end{proof}

\begin{remark}
Equation~\eqref{n4:eq:chart-determinant} is an identity over \(\Q\), not a
numerical fit. Modular computation only selects candidate row and column
indices; the determinant and common triangularization are reconstructed
exactly over \(\Q\).
\end{remark}

\subsection{Adding one quadratic direction}\label{n4:sec:extension}

\begin{lemma}\label{n4:lem:extension}
If \(q\in\Sym^2V\setminus E\), then
\[
  \rank\delta\bigl((E+\Span(q))\otimes V\bigr)\geq575.
\]
\end{lemma}

\begin{proof}
Define the linear map
\[
  \theta_q:V\longrightarrow
  \coker\bigl(\delta(E\otimes V)\bigr),\qquad
  v\longmapsto\psi_v([q]).
\]
Suppose independent \(v,w\) lie in \(\ker\theta_q\). Since \([q]\neq0\),
Proposition~\ref{n4:prop:chart} gives nonzero scalars \(a,b\) such that
\([q]=a[v^2]=b[w^2]\) in \(\Sym^2V/E\). Thus
\[
  av^2-bw^2\in E.
\]
The left side is nonzero and has rank at most two, contradicting
Theorem~\ref{n4:thm:rank-four-classification}. Hence
\(\dim\ker\theta_q\leq1\) and \(\rank\theta_q\geq15\). Its image is exactly
the new image contributed by \(q\otimes V\), so
\[
  \rank\delta\bigl((E+\Span(q))\otimes V\bigr)\geq560+15=575.
\]
\end{proof}

\subsection{A rank inequality}\label{n4:sec:double}

\begin{lemma}\label{n4:lem:double}
For two \(m\times n\) matrices \(A,B\),
\begin{equation}\label{n4:eq:double}
  \rank(A-B)\geq
  \rank[\,A\ B\,]
  +\rank\begin{bmatrix}A\\B\end{bmatrix}
  -\rank A-\rank B.
\end{equation}
\end{lemma}

\begin{proof}
Let
\[
  c=\dim(\im A\cap\im B),\qquad
  d=\dim(\row A\cap\row B).
\]
Then the horizontal and vertical concatenations have ranks
\(\rank A+\rank B-c\) and \(\rank A+\rank B-d\), respectively. For
\(x\in\ker(A-B)\), the vector \(Ax=Bx\) lies in \(\im A\cap\im B\), and
the kernel of \(x\mapsto Ax\) on \(\ker(A-B)\) is
\(\ker A\cap\ker B\). Therefore
\[
  \dim\ker(A-B)\leq c+n-\rank A-\rank B+d.
\]
Taking codimension in \(n\) gives \eqref{n4:eq:double}.
\end{proof}

\begin{remark}\label{n4:rem:scaling}
Replacing \(B\) by \(\alpha B\), \(\alpha\neq0\), does not change either
concatenation rank. Thus \eqref{n4:eq:double} also controls
\(A-\alpha B\).
\end{remark}

\subsection{Excluding seven terms}\label{n4:sec:lower}

\begin{theorem}\label{n4:thm:lower}
\[
  \ChowRank(\perm_4)\geq8.
\]
\end{theorem}

\begin{proof}
Fix a nonzero Chow term \(T\), and set
\[
  A=K(\perm_4),\qquad B=K(T),\qquad
  F=E_T,qquad d=\dim(E\cap F).
\]
Corollary~\ref{n4:cor:intersection} gives \(d\in\{0,1\}\).

If \(d=0\), then \eqref{n4:eq:row-containment} gives
\[
  \row A\subseteq E\otimes V^*,\qquad
  \row B\subseteq F\otimes V^*,
\]
and these ambient spaces intersect trivially. Thus the vertical concatenation
has rank \(\rank A+\rank B\). Lemma~\ref{n4:lem:double}, together with
\(\rank[\,A\ B\,]\geq\rank A=560\), gives
\begin{equation}\label{n4:eq:case0}
  \rank(A-\alpha B)\geq560\qquad(\alpha\neq0).
\end{equation}

If \(d=1\), choose \(0\neq q\in E\cap F\). Since
\(F\subseteq\Sym^2\Span(\ell_1,\ldots,\ell_4)\) and \(q\) has rank four,
the factors span a four-dimensional space and are independent. Hence
\(\dim F=6\) and \(\rank B=92\). Choose \(r\in F\setminus E\). The image
of the horizontal concatenation contains
\(\delta((E+\Span(r))\otimes V)\), so Lemma~\ref{n4:lem:extension} gives
\begin{equation}\label{n4:eq:horizontal}
  \rank[\,A\ B\,]\geq575.
\end{equation}
Also,
\[
  \dim(\row A\cap\row B)
  \leq\dim((E\cap F)\otimes V^*)=16,
\]
and hence
\begin{equation}\label{n4:eq:vertical}
  \rank\begin{bmatrix}A\\B\end{bmatrix}\geq560+92-16=636.
\end{equation}
Lemma~\ref{n4:lem:double} now gives
\begin{equation}\label{n4:eq:case1}
  \rank(A-\alpha B)\geq575+636-560-92=559
  \qquad(\alpha\neq0).
\end{equation}

Thus subtracting any nonzero Chow term from \(\perm_4\) leaves a flattening
of rank at least \(559\). If \(\perm_4\) had a seven-term decomposition,
remove one nonzero term. The remaining six terms would have flattening rank at
most \(6\cdot92=552\), contradicting \eqref{n4:eq:case0} or
\eqref{n4:eq:case1}. A decomposition with fewer nonzero terms gives the same
contradiction. Therefore \(\ChowRank(\perm_4)\geq8\).
\end{proof}

\subsection{The eight-term decomposition}\label{n4:sec:upper}

Glynn's identity is
\begin{equation}\label{n4:eq:glynn}
  \perm_n=2^{1-n}
  \sum_{\substack{\varepsilon\in\{\pm1\}^n\\\varepsilon_0=1}}
  \left(\prod_{i=0}^{n-1}\varepsilon_i\right)
  \prod_{j=0}^{n-1}
  \left(\sum_{i=0}^{n-1}\varepsilon_i x_{ij}\right).
\end{equation}
The sum has \(2^{n-1}\) terms, each a product of \(n\) linear forms. For
\(n=4\), it gives \(\ChowRank(\perm_4)\leq8\). Together with
Theorem~\ref{n4:thm:lower}, this proves Theorem~\ref{n4:thm:main}.

\subsection{Exact computation}\label{n4:sec:computation}

The computational part has a primary implementation and an independent
reconstruction.
\begingroup
\small
\begin{center}
\begin{tabularx}{0.94\textwidth}{@{}>{\hsize=0.78\hsize}L>{\hsize=1.22\hsize}L@{}}
\toprule
File & Purpose\\
\midrule
\path{perm4_quadratic_extension_chart_exact_verify.py}
& Reconstructs \eqref{n4:eq:659} over \(\Q\), computes the Schur complement
and \(\det S_0=-32768\), and checks simultaneous triangularization.\\
\path{perm4_quadratic_extension_full_minor_replay.py}
& Replays the full \(659\times659\) minor at several parameter points over
two primes as a separate cross-check.\\
\path{perm4_rank8_independent_verification.py}
& Rebuilds the ranks \(560\) and \(92\), the cubic prolongations, and the
chart minor directly from the displayed bases without importing project
modules.\\
\path{perm4_rank8_verify_all.py}
& Runs the upper bound, base ranks, chart calculation, and full-minor replay.\\
\bottomrule
\end{tabularx}
\end{center}

The independent program reports
\begin{verbatim}
independent_perm3_prolongations=16
independent_chow_koszul_rank=92
independent_chow_triple_prolongations=4
direct_E_image_rank=560
independent_chart_constant=-32768
independent_chart_triangular=True
INDEPENDENT_RANK8_AUDIT_PASS
\end{verbatim}
The lower bounds \(560\) and \(92\) come from displayed integer minors that
are nonzero modulo \(1000003\); the matching upper bounds come from the
sixteen and four explicit cubic prolongations. The identity
\(\det\mathcal M(v)=-32768v_{00}^{99}\) is verified over \(\Q\), and
\(S_4\times S_4\) symmetry carries that chart to every nonzero \(v\).

From the \texttt{perm4/certificates} directory, run
\begin{verbatim}
python perm4_rank8_independent_verification.py
python perm4_rank8_verify_all.py
\end{verbatim}
Both commands must exit successfully; the first prints
\texttt{INDEPENDENT\_RANK8\_AUDIT\_PASS}, and the second prints its aggregate
success marker.
\endgroup
